% Generated from the matching Typst scaffold by
% scripts/migrate_typst_report_to_latex.py.

% ── 7.1  Conclusion ──────────────────────────────────────────────────────
% PURPOSE: Restate the thesis and recap the metric → trade-off in one tight pass.
% MOVE: Restate + metric → trade-off recap (per the storytelling template).
% BEATS: see ordered bullets below — what we measured, the result, the reframe, the contribution.
% FIGURES/TABLES: none (cross-ref fig:pareto-footprint from Ch5/6).
% CITATIONS: none new (xiong2025 if interdependence is cited).
% SOURCE-DOC: README grounded-facts; docs/defense-qa.md (Q5).
% ─────────────────────────────────────────────────────────────────────────

\section{Conclusion}

\begin{todobox}
Expand the beats below into prose. RESTATE the study, recap the metric \textrightarrow{} trade-off, end by naming the contribution as a template, not a new policy. Thread "retrieval held constant", "non-inferiority", "footprint", "not every task needs memory". Numbers stay placeholders --- full 144-run matrix still RUNNING.
\end{todobox}

% - RESTATE: this thesis is a controlled measurement of six forgetting policies for AI
%   coding agents on SWE-Bench-CL, with retrieval held constant (pure cosine, identical
%   across all six) so the storage/forgetting policy is the only moving part.
% - METRIC → TRADE-OFF RECAP: on CL-F1, forgetting is non-inferior to full-memory
%   accumulation [RESULT: TOST, SESOI ±0.03] at materially lower memory footprint
%   [RESULT: …] — the footprint axis (A4) is where forgetting wins.
% - THE REFRAME: because only ~⅓ of tasks are interdependent [RESULT: ~0.30 overlap /
%   ~0.32 frac], memory is best understood as a DOSAGE / TARGETING question — "not every
%   task needs memory" — rather than an always-on accumulation default.
% - THE CONTRIBUTION: a clean, reproducible experimental TEMPLATE (retrieval-fixed
%   isolation, No-Memory↔Full-Memory bracketing, two-axis Pareto, interdependence
%   quantification, 144-run artifact) — NOT a new memory policy.
